\section{Known Effective Action}

We will now discuss 12d perspectives on the type IIB effective action, and argue that type IIB effective actions can not be written in 12d covariant form.
This is separate from the arguments related to the Eisenstein series structure, as well as the U(1)-violating sector, which has been recognized as difficult to be understood from 12d perspective \cite{Minasian:2015bxa,Liu:2025uqu}.
However, nevertheless, one may still hope that 12d contractions, upon reduction to 10d, and after applying LSZ reduction, produce the kinematic structure of string scattering amplitudes.
We argue that this can not be done.

The 4-point 8-derivative effective action can be written \cite{Minasian:2015bxa,Liu:2025uqu,Policastro:2008hg} as
\begin{equation}
\begin{aligned}
\frac{L_{\text{4-pt}}^{(3)}}{\alpha\, f_{0}(\tau,\bar{\tau})}
&= t_{8}t_{8} \biggl[R^4+6R^{2}\left(4|\delta\nabla P|^{2}+|\nabla G_{3}|^{2}\right)+24|\delta \nabla P|^{2}|\nabla G_{3}|^{2}\\
&+12R\left(\delta \nabla P(\nabla\overline{G}_{3})^{2}+D\overline{P}(\nabla{G}_{3})^{2}\right)\biggr]\\
&+1536 |\nabla_m P_n \nabla^m P^n|^2\\
&+O_{2}\left((|\nabla G_{3}|^{2})^{2}\right)
\label{eq:PTaction}
\end{aligned}
\end{equation}
where $\alpha = \frac{(\alpha')^3}{3 \cdot 2^{12}}$ and $(\delta \nabla P)_{mrns} \equiv g_{[m[n} \nabla_{r]} P_{s]}$ is defined to have the same symmetries as the Riemann tensor. 
Restricting to the gravity + scalar sector of \eqref{eq:PTaction}, we have
\begin{equation}
\frac{L_{\text{gravity+scalar}}}{\alpha f_0(\tau,\bar{\tau})} = t_8t_8
(R^4 + 24 R^2 |\delta \nabla P|^2) + 1536 (\nabla_m P_n \nabla^m P^n)(\nabla_r \bar{P}_s \nabla^r \bar{P}^s)
\label{eq:4pt-gravity-axiodilaton}
\end{equation}
Recall that $P_m \equiv \frac{i\nabla_m \tau}{2\tau_2}$, so schematically $\nabla P \sim \partial^2 \tau + (\partial \tau)^2$.
In the 12d perspective $\tau$ is embedded in the metric.
Suppose certain axio-dilaton couplings in the 10d effective action admits a 12d interpretation, one shall be able to rewrite such terms using 12d-covariant differential operators acting on the metric, i.e. the Riemann tensor. 

In particular, for the 12d metric ansatze we are working with, one may adopt complex tangent frame coordinates $z, \bar{z}$, with vielbein $e_{\alpha}{}^a = \frac{1}{2\sqrt{\tau_2}}\begin{pmatrix}
    -\tau & 1 \\ -\bar{\tau} & 1
\end{pmatrix}$, the nonvanishing components of the 12d Riemann tensor are \cite{Liu:2025uqu}
\begin{equation}
\hat{R}_{mrns} = R_{mrns},\quad
R_{z\bar{z}z\bar{z}} = -\frac14 |P|^2,\quad
R_{mnz\bar{z}} = P_{[m}\bar{P}_{n]},\quad
R_{mzn\bar{z}} = - \frac{P_m \bar{P}_n}{2}\quad
R_{mznz} = \frac{D_m P_n}{2}.
\end{equation}
% and equations of motions are
% \begin{equation}
% R_{mn} = 2 P_{(m} \bar{P}_{n)} \quad \quad D^m P_m = 0.
% \end{equation}
At the two derivative level, the 10d type IIB action can be written as
\begin{equation}
2\kappa^2 S_{IIB}
= \int d^{10}x \sqrt{-g} \left(R - 2 P_m \bar{P}^m\right).
\end{equation}
Substitute $P\bar{P}$ with $R_{mzn\bar{z}}$, we find
\begin{equation}
\sqrt{-g}[R - 2 P_m \bar{P}^m]
=\sqrt{-\hat{g}}[R + \hat{R}^{mz}{}_{mz} + \hat{R}^{m\bar{z}}{}_{m\bar{z}}]
=\sqrt{-\hat{g}} \hat{R},
\end{equation}
which is the usual story between the classical axio-dilaton sector and 12d gravity.

hence a term involving second derivatives of $\tau$ such as $\nabla P$.


% A general strategy for uplifting type IIB effective action may look like the following. Examine the effective action, substituting $\tau_1, \tau_2$ for components of the 12d metric on $T_2$, and derivatives of $\tau_1, \tau_2$ for 12d Riemannt tensor.
% This packages SL(2, Z) as part of 12d covariance. 
% One may then examine whether the resulting expression, in particular the coefficient of the distinct contraction patterns, may have a natural interpretation as arising from a higher dimensional contraction.
% The idea is that if certain terms of the 10d effective action arises from a 12d contraction, then one must be able to rewrite the 10d effective action to invert the index splitting, and promote a group of 10d terms to a 12d contraction.


% We note that crucially, while for a generic 12d metric ansatze, $R_{ambn}, R_{mnab}, R_{abcd}$ are independent of each other, for the specific F-theory ansatze, $R_{mnab}$ and $R_{abcd}$ can both be obtained as linear combinations or contractions of $R_{ambn}$.
% For example, $R_{z\bar{z}z\bar{z}} = \tfrac12 R^m{}_{zm\bar{z}}$, and $R_{mnz\bar{z}} = R_{m\bar{z}nz} - R_{mzn\bar{z}}$, leaving us with only 4 linearly independent components.
% This basis mismatch is related to the specific 12d metric ansatze chosen, and has important consequences as we will see.

The same story extends to the 4-point 8-derivative level.
Expanding the $t_8 t_8 R^2 |DP|^2$ sector, we find (with $\alpha = \zeta(3) \cdot 3 \cdot 2^8 / g_s^{3/2}$)
\begin{equation}
\begin{aligned}
\frac{L_{\text{grav}+\text{scalar}}}{\alpha}
&=  \frac{t_8 t_8 R^4}{3 \cdot 2^7}
+ 4\, D^{m}P^{n}\, D^{r}\bar{P}^{s}\, R_{m}{}^{t}{}_{r}{}^{u}\, R_{nust} \nonumber \\
&+ 4\, D^{m}P^{n}\, D^{r}\bar{P}^{s}\, R_{m}{}^{t}{}_{n}{}^{u}\, R_{rtsu}\\
&+ 4\, D^{m}P^{n}\, D_{m}P_{n}\, D^{r}\bar{P}^{s}\, D_{r}\bar{P}_{s} \nonumber \\
&- 4\, D_{m}\bar{P}^{r}\, D^{m}P^{n}\, R_{n}{}^{stu}\, R_{rstu} \\
&+ \frac{1}{2}\, D_{m}\bar{P}_{n}\, D^{m}P^{n}\, R^{rstu}\, R_{rstu} \nonumber + O(R^2 P^3).
\end{aligned}
\end{equation}
we note that because the only 12d Riemann tensor linear in $P$ is $R_{mznz}$ and its conjugate, we have
\begin{equation}
\hat{R}_{ambn} \hat{R}^{a}{}_{r}{}^{b}{}_{s} = 2 \text{Re} D_m P_n D_r \bar{P}_s + O(P^3).
\end{equation}
Recall that $D_m P_n$ is symmetric in $m,n$, substituting in the 12d Riemann tensors for coavariant derivatives on $P$, as well as recalling that $t_8 t_8 R^4$ can be written as
\begin{equation}
\begin{aligned}
\frac{t_8t_8 R^4}{3 \cdot 2^7}
&=  R_{pnqm} R^{prqs} R^{mt}{}_{ru} R^{nu}{}_{st} + \frac{1}{2}R_{mpnq} R^{qrps} R_{tsur} R^{untm} \\
& - \frac{1}{2}R_{mnrs} R^{rsnp} R_{pquv} R^{qmuv} - \frac{1}{4} R_{mnrs} R^{rs}{}_{pq} R^{qmuv} R_{uv}{}^{np}\\
& + \frac{1}{16} R_{mnrs} R^{mnpq} R_{pquv} R^{uvrs} + \frac{1}{32}(R_{mnrs} R^{mnrs})^2
\end{aligned}
\end{equation}
We find that
\begin{equation}
\begin{aligned}
\frac{L_{\text{grav}+\text{scalar}}}{\alpha}
&=
R_{pnqm} R^{prqs} R^{mt}{}_{ru} R^{nu}{}_{st}
+ 2 \hat{R}_{anbm} \hat{R}^{arbs} R^{mt}{}_{ru} R^{nu}{}_{st}\\
& + \frac{1}{2}\left[ R_{mpnq} R^{qrps} R_{tsur} R^{untm}
+ 4 \hat{R}_{ambn} \hat{R}^{bras} R_{tsur} R^{untm}
+ 2 \hat{R}_{ambn} \hat{R}^{bras} \hat{R}_{csdr} \hat{R}^{dncm}\right]\\
& - \frac{1}{2}\left[ R_{mnrs} R^{rsnp} R_{pquv} R^{qmuv}
 + 4 \hat{R}_{mban} \hat{R}^{anbp} R_{pquv} R^{qmuv} \right]\\
& - \frac{1}{4} R_{mnrs} R^{rs}{}_{pq} R^{qmuv} R_{uv}{}^{np}\\
& + \frac{1}{16} R_{mnrs} R^{mnpq} R_{pquv} R^{uvrs}\\
& + \frac{1}{32}(R_{rstu} R^{rstu})
\left[R_{mnrs} R^{mnrs}+ 8 \hat{R}_{ambn} \hat{R}^{ambn}\right]+ O(R^2 P^3)
\end{aligned}
\end{equation}

In this expression we have grouped terms with the same index contraction structure together. Because the only 12d Riemann tensor that is linear in $P$ is $R_{ambn}$, so we have, for example 
\begin{equation}
\hat{R}_{MNRS}\hat{R}^{MNRS} = R_{mrns} R^{mnrs} + 4 \hat{R}_{ambn} \hat{R}^{ambn} + O(P^3),
\end{equation}
hence the last term is compatible with coming from a 12d contraction $\frac{1}{32} (R_{MNRS} R^{MNRS})^2  \in t_8 t_8 \hat{R}^4$.
The rest of the terms also arise from 12d contraction $t_8 t_8 \hat{R}^4$ at 4pt, as was shown in \cite{Minasian:2015bxa}, using a ``top-down" approach, by making an educated guess of the 12d contraction.

The ``bottom-up" approach outlined here, despite requires more index-relabeling and tensor identities, should be the more systematic one. If there were to exist a 12d covariant expression that reproduce any 10d effective action, then after substituting 10d objects for 12d objects, one should always find that the terms group into a more natural pattern that is a result of index splitting. 
What would be a genuine obstruction, and this is independent of the human brain's capacity to identify a pattern, is if certain terms can not be written using 12d objects at all.

At 5-pt, the 

